\section{结论}
本文围绕动力电池包这一具体产品，对其智能装配与检测问题进行了系统分析。通过梳理产品构成、装配任务、工艺流程和检测链条，可以看出动力电池包制造具有结构约束强、电连接风险高、检测节点多和工艺耦合深等特点，因此适合作为智能装配与检测研究的典型对象。

在技术层面，本文重点分析了机器人、人机协作、机器视觉、AI、知识图谱和大模型智能体在装配过程中的应用逻辑。其中，机器人和协作机器人主要承担执行与辅助执行功能，机器视觉和深度学习方法负责感知与质量判断，知识图谱负责组织工艺知识和约束关系，大模型智能体则更适合作为问答、诊断和辅助决策工具，而非直接替代执行系统。这样的分层认识有助于避免将前沿技术简单堆砌为概念列表，也使多种技术能够在同一装配系统中形成分工明确的协同关系 \cite{Wang2022RCIM,Wang2025CJME,Zhang2023JMSYKG,Ji2024SciRep}。

在此基础上，本文提出了知识驱动的人机协同动力电池包智能装配与在线检测模式。该模式的核心并非追求完全无人化，而是通过知识图谱显式化工艺约束，通过视觉和传感器构建在线检测闭环，通过人机协作合理分配柔性任务与标准化任务，并利用大模型智能体增强现场知识调用与异常处置效率。相较于传统分散式自动化方案，这一模式更强调“工艺、数据、知识和人员”的统一协同，也更符合当前制造现场渐进式升级的现实路径 \cite{Sgarbossa2024JMSY,Long2025IJPR,Yao2023JMSY}。

总体而言，动力电池包智能装配的未来方向不是单点设备升级，而是系统级能力重构。只有真正打通装配过程、检测过程和知识决策过程，动力电池包产线才能在保证安全与质量的同时实现更高柔性、更强适应性和更好的可追溯性。本文的重点不在于给出唯一答案，而在于围绕动力电池包这一具体对象建立一套较完整的分析框架：以工艺分析为起点，以多技术融合为支撑，以新装配模式设计为落点，并进一步落实到系统构成、关键技术路径、应用前景和局限分析。

从面向具体产品的工程视角看，动力电池包智能装配研究最终仍需落回到模组层和整包层的真实结构关系上。图~\ref{fig:conclusion-module-structure}所示模组内部同时包含电芯堆叠、PE 隔离件、压力板、侧板和铝制汇流排等关键要素，这些细节不仅决定了后续装配顺序与检测位置，也直接影响自动化执行、质量控制和异常返修的边界。因此，智能装配的价值不在于脱离产品细节抽象讨论技术，而在于让技术能力稳定作用于具体产品结构和工艺约束 \cite{Kampker2021JRem,Baazouzi2022ApEnergy,Zhang2022JMSY,Wegener2022Energies}。

\begin{figure}[H]
  \centering
  \includegraphics[width=0.66\textwidth]{fig7_module_comparison.png}
  \caption{动力电池模组内部隔离、压紧与汇流排连接}
  \label{fig:conclusion-module-structure}
\end{figure}
